\section{Was der Träger für Materie und Eichstruktur leisten müsste}
\label{sec:shared-carrier}

\subsection{Die Biquaternionen als Kandidat}

Für die Feldrechnung in Abschnitt \ref{sec:frames} genügte ein komplexer Dreiervektor. Für einen gemeinsamen Träger suchen wir darüber hinaus eine Algebra, in der sich Zustände und Operationen zusammensetzen lassen. Der erste Test für einen solchen Kandidaten ist konkret. Er soll die Feldkoordinaten aufnehmen und ihre Lorentzwirkung samt der Invariante \(s\) wiedergeben. Für die physikalische Erweiterung muss er außerdem Raum für die inneren Symmetrien und die Zustandsauswahl lassen.

Der Companion bettet den Feldvektor in den Vektorteil eines Biquaternions ein. Mit den Quaternion-Einheiten \(e_k\) lautet der Ansatz
\begin{equation}
 q=it\,1+z_1e_1+z_2e_2+z_3e_3\in\mathbb C\otimes_{\mathbb R}\mathbb H,\qquad t\in\mathbb R .
 \label{eq:biquaternion}
\end{equation}
Für die quadratische Norm mit rein quaternionischer Konjugation gilt \(N(q)=-t^2+s(\mathbf z)\). Die Konjugationswirkung durch Biquaternionen der Norm eins trägt auf dem Vektorteil die eigentliche orthochrone Lorentzwirkung und erhält \(s\). Der Skalarparameter \(t\) bleibt unter dieser Wirkung fest. Seine mögliche Deutung ist eine zusätzliche Frage \cite{feldkern,grundidee}.

\paragraph{Welche Invariante der Träger tatsächlich erhält.}
Sei \(q=it\,1+\mathfrak z\) mit \(\mathfrak z=\sum_k z_k e_k\). Aus \(e_j e_k+e_ke_j=-2\delta_{jk}\,1\) folgt \(\mathfrak z^2=-s(\mathbf z)\,1\). Die quadratische Norm \(N(q)\) und ihre Multiplikativität lassen sich damit ausrechnen.
\begin{align}
 q\bar q&=(it\,1+\mathfrak z)(it\,1-\mathfrak z)
       =[-t^2+s(\mathbf z)]\,1,\\
 N(q_1q_2)&=N(q_1)N(q_2).
 \label{eq:biquaternion-check}
\end{align}
Die Produktidentität gilt für \(q_1,q_2\in\mathbb C\otimes_{\mathbb R}\mathbb H\) mit \(N(q)=q\bar q\) auf der vollen Algebra. Die Familie aus \eqref{eq:biquaternion} ist im Allgemeinen nicht unter Multiplikation abgeschlossen. Beispielsweise ist \(e_1^2=-1\) und hat reellen Skalarteil. Die Identität benutzt \(\overline{q_1q_2}=\bar q_2\bar q_1\) und die Zentralität der Norm. Für die Konjugation \(q\mapsto uqu^{-1}\) mit \(N(u)=1\) bleibt der Skalarteil \(it\,1\) fest. Über die Matrixrealisierung der Biquaternionen ist \(N\) die Determinante und die Norm-eins-Gruppe \(SL(2,\mathbb C)\). Ihre Konjugation auf dem spurfreien Vektorteil ergibt die Abbildung nach \(SO(3,\mathbb C)\), mit Kern \(\{\pm 1\}\) \cite{feldkern}. Auf diesem Teil können wir damit die Lorentz-Feldinvariante aus \eqref{eq:field-invariant} im selben algebraischen Träger behalten.

Dabei treten zwei Nullbedingungen auf. Für das Feld gilt \(s(\mathbf z)=0\), für den gesamten Kern \(N(q)=0\), also \(s(\mathbf z)=t^2\). Auf einem lichtartigen Feld ist \(N(q)=-t^2\). Der Parameter \(t\) bleibt in dieser Feldrepräsentation ein Lorentzskalar. Eine physikalische Deutung des gesamten Kerns muss angeben, welche der beiden Nullbedingungen sie verwendet und welche Rolle der Skalar dabei spielt \cite{feldkern,grundidee}.

\paragraph{Ein lichtartiges Feld in Spinorkoordinaten.}
Für \(\psi=(\psi_1,\psi_2)\in\mathbb C^2\) parametrisiert der folgende Ausdruck die lichtartige Feldklasse.
\begin{equation}
 \mathbf z(\psi)=
 (\psi_1^2-\psi_2^2,\,i(\psi_1^2+\psi_2^2),\,-2\psi_1\psi_2).
 \label{eq:field-null-spinor}
\end{equation}
Die beiden ersten Quadrate ergeben zusammen \(-4\psi_1^2\psi_2^2\), das dritte gibt den Gegenbeitrag. Die Welle aus \eqref{eq:field-example} entspricht \(\psi=(\sqrt a,0)\) in den dort verwendeten Feldkoordinaten. Jedes von null verschiedene \(\mathbf z\) mit \(s(\mathbf z)=0\) hat die beiden Urbilder \(\pm\psi\). Man gewinnt sie aus
\[
 \psi_1^2=\tfrac12(z_1-iz_2),\qquad
 \psi_2^2=-\tfrac12(z_1+iz_2),\qquad
 \psi_1\psi_2=-\tfrac12 z_3 .
\]
Das ist eine algebraische Beschreibung der klassischen Feldklasse. Für ein quantisiertes Einphotonensystem wären Zustand und Observablen zusätzlich zu spezifizieren \cite{feldkern}.

So erhält der Ansatz einen konkreten algebraischen Kandidaten. Auf seinem Feldteil lassen sich die Feldklassen und die reell-lineare Austauschoperation beschreiben. Für die nächste Aufgabe brauchen wir einen Zustandsraum mit einem Skalarprodukt und einer erklärten Auswahl der Materiezustände. Der folgende Vergleich untersucht, wie eine solche Auswahl drei geeignete Zustände liefern könnte. Welche physikalische Darstellung sie realisiert, muss eigens gezeigt werden.

\subsection{Materiezustände und Generationen}
\label{sec:matter}

Ein Generationsindex soll Zustände unterscheiden, die in den relevanten inneren Ladungen übereinstimmen. Aus einem dreidimensionalen Unterraum allein folgt daher zunächst nur die Dimension. Ob die drei Richtungen Farben, Generationen oder eine andere Multiplizität bezeichnen, entscheidet die Wirkung der inneren Symmetrien und der Observablen auf diesen Zuständen.

Eine Vergleichskonstruktion des Companion macht die Wahl des Unterraums explizit \cite{companion,physics}. Sei \(T_0:\mathbb C^3\to\mathcal K\) eine Isometrie in einen komplexen Hilbertraum und \(P_c\) der orthogonale Projektor auf die im Vergleichsmodell zulässigen Zustände. Die Matrix \(A_0=T_0^*P_cT_0\) erfüllt
\[
 \langle v,A_0v\rangle=\|P_cT_0v\|^2 .
\]
Die Bedingung \(A_0>0\) besagt, dass die drei projizierten Ausgangszustände linear unabhängig bleiben. Dann ist
\begin{equation}
 S=P_cT_0A_0^{-1/2}, \qquad
 S^*S=A_0^{-1/2}T_0^*P_c^2T_0A_0^{-1/2}=\id_3,\qquad P_cS=S .
 \label{eq:seed}
\end{equation}
Damit erhält man drei orthonormale Zustände im Bild von \(P_c\). Die Wahl von \(T_0\) und \(P_c\) geht als Voraussetzung ein. \(P_c\) setzt die Zulässigkeit im Vergleichsraum um, \(C\) aus \eqref{eq:projections} die HC-Gradierung.

Ein Operator \(O\) auf \(\mathcal K\) lässt sich zu \(S^*OS\) komprimieren, sofern das Bild von \(S\) in seinem Definitionsbereich liegt. Erhält \(O\) dieses Bild, gilt \(OS=S(S^*OS)\). Die Kompression stellt dann die Einschränkung des Operators dar. Im allgemeinen Fall erfasst sie die Matrixelemente zwischen den drei ausgewählten Zuständen, während Übergänge aus ihrem Unterraum heraus zusätzlich zu untersuchen sind. Das betrifft insbesondere die Zuordnung eines Massenspektrums zu einem nativen Operator.

\paragraph{Eine positive Matrix aus der gewählten Zustandsauswahl.}
Der Companion untersucht eine Abbildung aus dem ausgewählten Zustandsraum in dessen orthogonales Komplement. Sei \(\widehat R\) ein beschränkter, dimensionslos normierter Operator auf \(\mathcal K\). Auf den Zuständen \(\psi_i=S\mathbf e_i\) setzt der Vergleich \(B_\perp=(I-P_c)\widehat R S\) an. Die daraus gebildete Matrix ist
\begin{equation}
 Q_\perp=B_\perp^*B_\perp,\qquad
 (Q_\perp)_{ij}
 =\big\langle (I-P_c)\widehat R\psi_i,(I-P_c)\widehat R\psi_j\big\rangle .
 \label{eq:generation-defect-readout}
\end{equation}
Für \(v\in\mathbb C^3\) ist \(\langle v,Q_\perp v\rangle=\|B_\perp v\|^2\). Außerdem gilt \([\widehat R,P_c]\psi_i=(I-P_c)\widehat R\psi_i\). Dessen Erwartungswert in \(\psi_i\) verschwindet wegen der Orthogonalität zum Bild von \(P_c\), auch wenn die Norm des Vektors positiv ist. Eine Auslese durch den Erwartungswert und eine Auslese durch die Übergangsnorm liefern damit unterschiedliche Größen. Bei exakter Kommutation \([\widehat R,P_c]=0\) verschwindet auch die gesamte Matrix \(Q_\perp\) \cite{companion,repo}.

Welche Zahlen würde man bei \(Q_\perp\ne0\) als Massenkandidaten lesen? Setzt man im Vergleich eine Skala \(\kappa>0\) und die effektive Matrix \(\kappa Q_\perp\) an, sind dafür ihre Eigenwerte relevant. Die Diagonale in den ursprünglichen Generationenkoordinaten enthält die drei quadrierten Normen der einzelnen Übergänge. Wie weit beides auseinanderliegen kann, zeigt die im Companion angegebene Rang-eins-Matrix \(Q=\mathbf1\mathbf1^*/3\) mit \(\mathbf1=(1,1,1)^{\mathsf T}\). Ihre Diagonale ist \((1/3,1/3,1/3)\). Aus \(Q\mathbf1=\mathbf1\) und \(Q|_{\mathbf1^\perp}=0\) folgen die Eigenwerte \((1,0,0)\). In der Koide-Kennzahl aus Abschnitt \ref{sec:mass} liefern diese beiden Auswertungen folglich \(1/3\) und \(1\). Der Vergleich macht sichtbar, welche Ausleseregel, welcher Operator und welche Skala feststehen müssen, bevor ein Massenspektrum bestimmt werden kann \cite{companion,repo}.

Eine physikalische Generationskonstruktion müsste aus dem gemeinsamen Träger den Ausgangsraum auswählen, die Ladungsdarstellung auf ihm bestimmen und zeigen, wie Massen die drei Zustände unterscheiden. Gleichung \eqref{eq:seed} legt dabei offen, an welcher Stelle die Wahl der Zustände eingeht. Der Koide-Test in Abschnitt \ref{sec:mass} setzt daran an, sobald ein Kandidat für die Massenauslese feststeht.
